% =====================================================================
% HOJA 8 (viii) – RESUMEN Y PALABRAS CLAVES
% =====================================================================
\begin{center}
    \vspace*{0.5cm}
    {\bfseries RESUMEN}
\end{center}
\vspace{0.8cm}

\justifying
El presente proyecto tiene como objetivo desarrollar ``Kiddy Quiz'', una aplicación web interactiva orientada a la evaluación y el refuerzo de matemáticas para estudiantes de educación básica elemental. Para la fundamentación y levantamiento de requisitos, se llevó a cabo una investigación de campo que incluyó entrevistas a 5 profesionales de la educación, 2 psicopedagogos, lo cual permitió identificar las barreras cognitivas y de ansiedad presentes en los procesos evaluativos tradicionales. El desarrollo del software se gestionó bajo la metodología ágil de Programación Extrema (XP), asegurando la calidad y robustez del código. La arquitectura del sistema se construyó utilizando Angular para el desarrollo del frontend, NestJS para la lógica del backend y PostgreSQL como motor de base de datos relacional. Como principal innovación, la plataforma integra el uso de interfaces basadas en pictogramas para facilitar la navegación autónoma infantil, así como la implementación de Inteligencia Artificial (Gemini) para proveer retroalimentación motivacional tras cada evaluación. La validación del sistema se realizó a través de la Escala de Usabilidad del Sistema (SUS) y el Modelo de Aceptación Tecnológica (TAM), demostrando una usabilidad en categoría ``Excelente'' y una intención de uso casi unánime por parte de los estudiantes. Además, para garantizar la protección de los datos y el correcto control de roles, se implementaron medidas de seguridad mediante autenticación con tokens JWT. Como resultado, se ha logrado consolidar una herramienta tecnológica inclusiva que optimiza la gestión académica del docente, reduce la frustración del estudiante y representa un aporte significativo en la modernización de la evaluación educativa infantil.

\vspace{0.5cm}

\noindent\textbf{Palabras clave:} aplicación web, educación básica, matemáticas, pictogramas, inteligencia artificial, usabilidad, Kiddy Quiz.

\newpage
